\input{../tex-inputs/latex-korrekturansicht-vorspann}

\section[Berta Zuckerkandl an Arthur Schnitzler, {[}29. 3. 1897?{]}]{L04015 Berta Zuckerkandl an Arthur Schnitzler, {[}29. 3. 1897?{]}}
\nopagebreak\mylabel{L04015v}
\rehead{ }\normalsize\beginnumbering\briefempfaengerindex{Schnitzler, Arthur@\textsc{Schnitzler, Arthur}!zzzZuckerkandl, Berta@\emph{von Berta Zuckerkandl}!1897-03-291@{{[}29. 3. 1897?{]}}|(be}
\toendnotes[C]{\smallbreak\pagebreak[2]}
\correspDesc{Versand  durch Berta Zuckerkandl am [29. 3. 1897?] in Wien
\newline{}Erhalt  durch Arthur Schnitzler im Zeitraum [29. 3. 1897 – 1. 4. 1897?] in Wien}\toendnotes[C]{\smallbreak}
\Standort{CUL, Schnitzler, B 200.}
\physDesc{Brief, 1 Blatt, 3 Seiten, 538 Zeichen
\newline{}Handschrift: schwarze Tinte, lateinische Kurrent
\newline{}Schnitzler: beschriftet: »Zuckerkandl« und datiert »97« }\toendnotes[C]{\smallbreak}
\pstart{}{\pb}Hochverehrter Herr Doktor!\pend\vspace{0.5em}
\pstart
           Erlauben sie mir Ihnen meinen herzlichsten, innigsten Dank für Ihre wirklich rührende
                  \label{K_L04015-1v}\edtext{Güte}{\lemma{\textnormal{\emph{Güte}}}\Cendnote{\textnormal{Das Korrespondenzstück
                  ist von der Verfasserin undatiert geblieben. \textcolor{blue}{Schnitzler} datiert es im Jahr »{[}18{]}97«.
                     Am 28. 3. 1897 trug 
                     \textcolor{blue}{Schnitzler} den 1. Akt von
                     \emph{\textcolor{green}{Freiwild}\pwindex{Schnitzler, Arthur 15. 5. 1862 Wien – 21. 10. 1931 ebd.@\textsc{Schnitzler, Arthur} (15. 5. 1862 Wien – 21. 10. 1931 ebd.), \emph{Schriftsteller, Mediziner}!Freiwild. Schauspiel in 3 Akten@\strich\emph{Freiwild. Schauspiel in 3 Akten}|pwk}} bei
                      einer gemeinsamen \textcolor{violet}{Lesung}\eventindex{Bösendorfer-Saal@\textbf{Bösendorfer-Saal}!Lesung von Hirschfeld, Hofmannsthal, Schnitzler und Bahr, 28.3.1897@Lesung von Hirschfeld, Hofmannsthal, Schnitzler und Bahr, 28.3.1897|pwk} mit \textcolor{blue}{Georg Hirschfeld}\pwindex{Hirschfeld, Georg 11.\,2.\,1873 Berlin – 17.\,1.\,1942 München@\textsc{Hirschfeld, Georg} (11.\,2.\,1873 Berlin – 17.\,1.\,1942 München), \emph{Schriftsteller}|pwk},
                     \textcolor{blue}{Hugo von Hofmannsthal}\pwindex{Hofmannsthal, Hugo von 1.\,2.\,1874 Wien – 15.\,7.\,1929 Rodaun@\textsc{Hofmannsthal, Hugo von} (1.\,2.\,1874 Wien – 15.\,7.\,1929 Rodaun), \emph{Schriftsteller}|pwk} und \textcolor{blue}{Hermann Bahr}\pwindex{Bahr, Hermann 19.\,7.\,1863 Linz – 15.\,1.\,1934 München@\textsc{Bahr, Hermann} (19.\,7.\,1863 Linz – 15.\,1.\,1934 München), \emph{Schriftsteller, Kritiker}|pwk} vor.
                     Das geschah für einen wohltätigen Zweck und wurde von \textcolor{blue}{Berta Zuckerkandl}\pwindex{Zuckerkandl, Berta 13.\,4.\,1864 Wien – 16.\,10.\,1945 Paris@\textsc{Zuckerkandl, Berta} (13.\,4.\,1864 Wien – 16.\,10.\,1945 Paris), \emph{Schriftstellerin, Journalistin, Übersetzerin}|pwk}
                     organisiert, wie das \emph{\textcolor{green}{Illustrirte Wiener Extrablatt}\pwindex{Illustrirtes Wiener Extrablatt@\emph{Illustrirtes Wiener Extrablatt}|pwk}} implizierte: »\textcolor{green}{Die geiſtreiche
                           \textcolor{blue}{Frau}\pwindex{Zuckerkandl, Berta 13.\,4.\,1864 Wien – 16.\,10.\,1945 Paris@\textsc{Zuckerkandl, Berta} (13.\,4.\,1864 Wien – 16.\,10.\,1945 Paris), \emph{Schriftstellerin, Journalistin, Übersetzerin}|pwv} eines berühmten \textcolor{blue}{Anatomen}\pwindex{Zuckerkandl, Emil 1.\,9.\,1849 Győr – 28.\,5.\,1910 Wien@\textsc{Zuckerkandl, Emil} (1.\,9.\,1849 Győr – 28.\,5.\,1910 Wien), \emph{Anatom}|pwv} kam auf die gute Idee, die vier Parnaßgänger für einen wohltätigen 
                     Zweck als Vorleser zu gewinnen.}\pwindex{Jung-Wien bei Bösendorfer@\emph{Jung-Wien bei Bösendorfer}|pwv}« ([O. V.]: \emph{\textcolor{green}{Jung-Wien bei Bösendorfer}\pwindex{Jung-Wien bei Bösendorfer@\emph{Jung-Wien bei Bösendorfer}|pwk}}. In: \emph{\textcolor{green}{Illustrirtes Wiener Extrablatt}\pwindex{Illustrirtes Wiener Extrablatt@\emph{Illustrirtes Wiener Extrablatt}|pwk}}, Jg. 26, Nr. 88, 29. 3. 1899, S. XXXX)
                  Beim vorliegenden Brief handelt es sich dementsprechend um das Dankschreiben, das am Folgetag verfasst sein dürfte.}}}\label{K_L04015-1} auszusprechen. Sie haben
               mit der edelsten Bereitwilligkeit Ihr grosses Talent{[},{]} Ihre grosse Zugkraft uns zur
               Verfügung gestellt. Ich bin wirklich ganz be{\pb}schämt, und hoffe nun dass der \uline{kolossale} Erfolg den ihr geradezu verblüffend
               gescheuter \textcolor{violet}{\textcolor{green}{Akt}\pwindex{Schnitzler, Arthur 15. 5. 1862 Wien – 21. 10. 1931 ebd.@\textsc{Schnitzler, Arthur} (15. 5. 1862 Wien – 21. 10. 1931 ebd.), \emph{Schriftsteller, Mediziner}!Freiwild. Schauspiel in 3 Akten@\strich\emph{Freiwild. Schauspiel in 3 Akten}|pwv}{}\ledrightnote{{$\rightarrow$}\emph{\textcolor{green}{Freiwild. Schauspiel in 3 Akten}}}}\eventindex{Urania Theater@\textbf{Urania Theater}!Premiere von Der Schleier der Pierrette, 16.3.1912@Premiere von Der Schleier der Pierrette, 16.3.1912|pwuv}{}\ledrightnote{{$\rightarrow$}\emph{\textcolor{violet}{Premiere von Der Schleier der Pierrette, 16.3.1912}}} errungen hat, Sie einigermaassen für die grosse Mühe
               entschädigt hat. Nehmen Sie geehrter Herr die Versicherung meiner tiefsten
               Dankbarkeit entgegen.\pend
           
\pstart
           {\pb}Ihre sehr ergebene {\\[\baselineskip]}\spacefill\mbox{Berta Zuckerkandl}\pend
           \leftskip=0em{}\selectlanguage{ngerman}\endnumbering\briefempfaengerindex{Schnitzler, Arthur@\textsc{Schnitzler, Arthur}!zzzZuckerkandl, Berta@\emph{von Berta Zuckerkandl}!1897-03-291@{{[}29. 3. 1897?{]}}|)be}\mylabel{L04015h}  \newcommand{\dateiname}{L04015}\newcommand{\titel}{Berta Zuckerkandl an Arthur Schnitzler, [29. 3. 1897?]}\newcommand{\editorInnen}{Herausgegeben von Jahnke, SelmaMüller, Martin Anton}\input{../tex-inputs/latex-korrekturansicht-abspann}
